\section{Discussion}\label{sec:discussion}

The results show that the relative accuracy of conditional-law learners depends on both the data-generating process and the causal estimand. WCF performs well in the location and shape-effect designs with assignment depending on prognostic covariates. The additional nonlinear designs provide a broader comparison, including settings in which a forest baseline is more accurate. Covariate-dependent assignment creates a need for adjustment, but does not alone determine which estimator will perform best. A constant treatment probability different from one half can describe a randomized experiment with unequal allocation; a nonconstant propensity can also depend on covariates unrelated to the outcome.

WCF combines two distinct forms of adjustment. Shared tree partitions and contrast shrinkage pool outcome information across treatment arms. Functional calibration uses estimated treatment probabilities to correct the reported averages of selected outcome functionals. The latter has the augmented-score robustness property of Section~\ref{sec:dr}. Neither adjustment eliminates the need for overlap, and the comparison with DRF and Causal-DRF does not isolate the contribution of calibration from all other architectural differences. Doubly robust distributional procedures already exist \citep{martineztaboada2023efficient,GRFDistributionalCausalEffects}; the contribution here is the combination of a Wasserstein particle representation, shared arm updates, and functional calibration, together with the reference-distance estimands.

The reference-distance effects express a directional causal question: does treatment reduce the expected distance from a unit's outcome distribution to a chosen benchmark? Their interpretation depends on the reference. A negative effect indicates average movement toward that benchmark, without implying improvement for every unit or every feature of the distribution. Substantive applications should justify the benchmark and consider alternative scientifically meaningful references.

The optional two-part construction applies to an atom at an entirely degenerate outcome distribution. It estimates the probability of that component separately from the remaining conditional law. The experiments evaluate this specific mixture structure. A dataset containing individual zeros within otherwise nondegenerate unit-level distributions presents a different representation problem.

Several limitations delimit these findings. The studies assume independently sampled units, observed quantile vectors, consistency, conditional exchangeability, and positivity. Applications based on estimated quantile curves or repeated observations within regions require additional treatment of measurement error and dependence. The reported fits use one grid resolution and particle budget; stability to these choices remains to be established. A fixed number of equal-weight particles limits the laws and functional means that can be represented, and propensity clipping can introduce bias at the extremes of the assignment distribution. The simulations do not assess confidence-interval coverage or establish nuisance convergence rates for the fitted learner. All empirical comparisons are synthetic; an applied evaluation remains necessary to assess the method under substantive identification assumptions.

\section{Conclusion}\label{sec:conclusion}

Wasserstein Causal Forests estimate conditional laws of distribution-valued outcomes and calibrate population and moderator-specific effects on their functionals. The reference-distance TATE and TCATE measure whether treatment moves these outcomes toward a prespecified distribution. Simulations demonstrate substantial gains over the forest comparators in several designs with observed confounding, together with variation across nonlinear designs and estimands. An explicit two-part representation extends the law estimator to a structural zero component. These results motivate further evaluation of WCF with alternative representation budgets, flexible propensity models, and applied distribution-valued outcomes.
